% GENERATED FILE. Edit canonical lessons, then run npm run generate:books.
% canonical-source-hash: fnv1a64:0776da5eb5152cc9
% canonical-lessons: SA-C288-aparadhah, SA-C288-sandehah, SA-C288-abhiyogah, SA-C288-tarjayati, SA-C288-sastram

\chapter{At the Watch-house}
\label{ch:sa-at-the-watch-house}
\hlchaptermodality{288}

\begin{chapteropening}
\textbf{By the end of this chapter:} \emph{I can report a crime and say what the charge depends on.}

Complete the last lesson of chapter 288: I can report a crime and say what the charge depends on.
\end{chapteropening}

\section[\texorpdfstring{\sk{अपराधः} (aparādhaḥ)}{अपराधः (aparādhaḥ)}]{\sk{अपराधः} (aparādhaḥ) — an offence; a crime}

\label{lesson:SA-C288-aparadhah}



\begin{quote}
Before the new one: say the Sanskrit for stays awake, then the Sanskrit for an uproar.
\end{quote}

\subsection*{You'll want to know: \sk{अपराधः}}
\textbf{\sk{अपराधः}} — \emph{aparādhaḥ} — \textquotedblleft{}an offence; a crime\textquotedblright{}.

\begin{quote}
\sk{चौर्यम्} \sk{अपराधः} \sk{अस्ति।}
\end{quote}

\begin{quote}
\emph{cauryam aparādhaḥ asti.}
\end{quote}

\begin{quote}
Theft is an offence.
\end{quote}

A masculine noun.

\begin{cousinweb}
\textbf{\sk{अपराधः}} (\emph{aparādhaḥ}) is built from \emph{apa}, \textquotedblleft{}away\textquotedblright{}, and the root √rādh, \textquotedblleft{}to succeed\textquotedblright{}: MW gives \emph{apa-rādh} as to miss one's aim and to offend, and \emph{aparādha} as an offence (MW).
\end{cousinweb}

\begin{grammarlens}[title={What you've built}]
One more word for telling what happened.
\end{grammarlens}

\subsection*{Your turn}
\begin{itemize}
\raggedright
  \item \emph{Say it:} \emph{aparādhaḥ}
  \item \emph{Say it:} \emph{aparādhaḥ}, once more
  \item \emph{Say it:} \emph{cauryam aparādhaḥ asti.} — Theft is an offence.
  \item \emph{Recall:} say the Sanskrit for stays awake, then the Sanskrit for an uproar, then say \emph{aparādhaḥ} again
\end{itemize}

\begin{tcolorbox}[breakable,colback=teal!4,colframe=teal!35!black,title={Before you move on}]
What does \sk{अपराधः} mean? (\textquotedblleft{}An offence; a crime\textquotedblright{}.) Say it once more.
\end{tcolorbox}
\section[\texorpdfstring{\sk{सन्देहः} (sandehaḥ)}{सन्देहः (sandehaḥ)}]{\sk{सन्देहः} (sandehaḥ) — a suspicion; a doubt}

\label{lesson:SA-C288-sandehah}



\begin{quote}
Before the new one: say the Sanskrit for an uproar, then the Sanskrit for an offence.
\end{quote}

\subsection*{You'll want to know: \sk{सन्देहः}}
\textbf{\sk{सन्देहः}} — \emph{sandehaḥ} — \textquotedblleft{}a suspicion; a doubt\textquotedblright{}.

\begin{quote}
\sk{मम} \sk{सन्देहः} \sk{अस्ति।}
\end{quote}

\begin{quote}
\emph{mama sandehaḥ asti.}
\end{quote}

\begin{quote}
I have a suspicion.
\end{quote}

A masculine noun.

\begin{cousinweb}
\textbf{\sk{सन्देहः}} (\emph{sandehaḥ}) is built from \emph{sam}, \textquotedblleft{}together\textquotedblright{}, and the root √dih, \textquotedblleft{}to smear, heap up\textquotedblright{}; MW's first sense is a mass stuck together, and its usual sense doubt or uncertainty, and a suspicion is its everyday modern use (MW).
\end{cousinweb}

\begin{grammarlens}[title={What you've built}]
One more word for telling what happened.
\end{grammarlens}

\subsection*{Your turn}
\begin{itemize}
\raggedright
  \item \emph{Say it:} \emph{sandehaḥ}
  \item \emph{Say it:} \emph{sandehaḥ}, once more
  \item \emph{Say it:} \emph{mama sandehaḥ asti.} — I have a suspicion.
  \item \emph{Recall:} say the Sanskrit for an uproar, then the Sanskrit for an offence, then say \emph{sandehaḥ} again
\end{itemize}

\begin{tcolorbox}[breakable,colback=teal!4,colframe=teal!35!black,title={Before you move on}]
What does \sk{सन्देहः} mean? (\textquotedblleft{}A suspicion; a doubt\textquotedblright{}.) Say it once more.
\end{tcolorbox}
\section[\texorpdfstring{\sk{अभियोगः} (abhiyogaḥ)}{अभियोगः (abhiyogaḥ)}]{\sk{अभियोगः} (abhiyogaḥ) — a charge; an accusation}

\label{lesson:SA-C288-abhiyogah}



\begin{quote}
Before the new one: say the Sanskrit for an offence, then the Sanskrit for a suspicion.
\end{quote}

\subsection*{You'll want to know: \sk{अभियोगः}}
\textbf{\sk{अभियोगः}} — \emph{abhiyogaḥ} — \textquotedblleft{}a charge; an accusation\textquotedblright{}.

\begin{quote}
\sk{अभियोगः} \sk{केवलं} \sk{सन्देहाधीनः।}
\end{quote}

\begin{quote}
\emph{abhiyogaḥ kevalaṃ sandehādhīnaḥ.}
\end{quote}

\begin{quote}
The charge rests only on suspicion.
\end{quote}

A masculine noun.

\begin{cousinweb}
\textbf{\sk{अभियोगः}} (\emph{abhiyogaḥ}) is built from \emph{abhi}, \textquotedblleft{}towards, against\textquotedblright{}, and the root √yuj, \textquotedblleft{}to join\textquotedblright{}; MW gives \emph{abhiyoga} as an effort, an attack and, in law, a charge or accusation (MW).
\end{cousinweb}

\begin{grammarlens}[title={What you've built}]
One more word for telling what happened.
\end{grammarlens}

\subsection*{Your turn}
\begin{itemize}
\raggedright
  \item \emph{Say it:} \emph{abhiyogaḥ}
  \item \emph{Say it:} \emph{abhiyogaḥ}, once more
  \item \emph{Say it:} \emph{abhiyogaḥ kevalaṃ sandehādhīnaḥ.} — The charge rests only on suspicion.
  \item \emph{Recall:} say the Sanskrit for an offence, then the Sanskrit for a suspicion, then say \emph{abhiyogaḥ} again
\end{itemize}

\begin{tcolorbox}[breakable,colback=teal!4,colframe=teal!35!black,title={Before you move on}]
What does \sk{अभियोगः} mean? (\textquotedblleft{}A charge; an accusation\textquotedblright{}.) Say it once more.
\end{tcolorbox}
\section[\texorpdfstring{\sk{तर्जयति} (tarjayati)}{तर्जयति (tarjayati)}]{\sk{तर्जयति} (tarjayati) — threatens}

\label{lesson:SA-C288-tarjayati}



\begin{quote}
Before the new one: say the Sanskrit for a suspicion, then the Sanskrit for a charge in law.
\end{quote}

\subsection*{You'll want to know: \sk{तर्जयति}}
\textbf{\sk{तर्जयति}} — \emph{tarjayati} — \textquotedblleft{}threatens\textquotedblright{}.

\begin{quote}
\sk{चौरः} \sk{बालं} \sk{तर्जयति।}
\end{quote}

\begin{quote}
\emph{cauraḥ bālaṃ tarjayati.}
\end{quote}

\begin{quote}
The thief threatens the boy.
\end{quote}

Past (he …): \sk{सः} \sk{बालं} \sk{तर्जितवान्}, \emph{saḥ bālaṃ tarjitavān}.

\begin{cousinweb}
\textbf{\sk{तर्जयति}} (\emph{tarjayati}) comes from the root √tarj, \textquotedblleft{}to threaten, scold\textquotedblright{}; \emph{tarjayati} is its causative, which MW gives with the same sense (MW).
\end{cousinweb}

\begin{grammarlens}[title={What you've built}]
One more word for telling what happened.
\end{grammarlens}

\subsection*{Your turn}
\begin{itemize}
\raggedright
  \item \emph{Say it:} \emph{tarjayati}
  \item \emph{Say it:} \emph{tarjayati}, once more
  \item \emph{Say it:} \emph{cauraḥ bālaṃ tarjayati.} — The thief threatens the boy.
  \item \emph{Recall:} say the Sanskrit for a suspicion, then the Sanskrit for a charge in law, then say \emph{tarjayati} again
\end{itemize}

\begin{tcolorbox}[breakable,colback=teal!4,colframe=teal!35!black,title={Before you move on}]
What does \sk{तर्जयति} mean? (\textquotedblleft{}Threatens\textquotedblright{}.) Say it once more.
\end{tcolorbox}
\section[\texorpdfstring{\sk{शस्त्रम्} (śastram)}{शस्त्रम् (śastram)}]{\sk{शस्त्रम्} (śastram) — a weapon}

\label{lesson:SA-C288-sastram}



\begin{quote}
Before the new one: say the Sanskrit for a charge in law, then the Sanskrit for threatens.
\end{quote}

\subsection*{You'll want to know: \sk{शस्त्रम्}}
\textbf{\sk{शस्त्रम्}} — \emph{śastram} — \textquotedblleft{}a weapon\textquotedblright{}.

\begin{quote}
\sk{चौरस्य} \sk{हस्ते} \sk{शस्त्रम्} \sk{अस्ति।}
\end{quote}

\begin{quote}
\emph{caurasya haste śastram asti.}
\end{quote}

\begin{quote}
The thief has a weapon in his hand.
\end{quote}

A neuter noun.

\begin{cousinweb}
\textbf{\sk{शस्त्रम्}} (\emph{śastram}) comes from the root √śas, \textquotedblleft{}to cut, wound\textquotedblright{}; MW gives \emph{śastra} as an instrument for cutting, a knife, a sword or any weapon (MW).
\end{cousinweb}

\begin{grammarlens}[title={What you've built}]
One more word for telling what happened.
\end{grammarlens}

\subsection*{Your turn}
\begin{itemize}
\raggedright
  \item \emph{Say it:} \emph{śastram}
  \item \emph{Say it:} \emph{śastram}, once more
  \item \emph{Say it:} \emph{caurasya haste śastram asti.} — The thief has a weapon in his hand.
  \item \emph{Recall:} say the Sanskrit for a charge in law, then the Sanskrit for threatens, then say \emph{śastram} again
\end{itemize}

\begin{tcolorbox}[breakable,colback=teal!4,colframe=teal!35!black,title={Before you move on}]
What does \sk{शस्त्रम्} mean? (\textquotedblleft{}A weapon\textquotedblright{}.) Say it once more.
\end{tcolorbox}
